\chapter{Evaluation Methodology}
\label{ch:evaluation_methodology}
This chapter presents the evaluation methodology applied to both unconditional and
conditional generated data. Multiple metrics are used to assess distributional similarity,
moment reproduction, path-wise behaviour, and conditional predictive quality, combining
quantitative metrics with visual insights.

\section{Distributional Similarity}

To assess the similarity between the real and generated data distributions, two
complementary metrics are reported: the Wasserstein distance \cite{arjovsky2017wasserstein},
which captures global distributional shifts, and the Kolmogorov--Smirnov (KS) test
\cite{massey1951kolmogorov}, which provides a nonparametric hypothesis test for
distributional differences. Both metrics are computed on the log-return datasets
\(\mathcal{D}_{\text{train}}\) and \(\mathcal{D}_{\text{gen}}\).

\textbf{Kolmogorov--Smirnov} \quad The KS statistic measures the maximum distance between
the empirical cumulative distribution functions of two samples. Given two empirical
cumulative distributions \(F(x)\) and \(G(x)\), it is defined as:
\[
    D_{\text{KS}} = \sup_x |F(x) - G(x)|
\]
where \(\sup_x\) denotes the supremum over all values of \(x\). The KS test is
advantageously nonparametric and sensitive to both location and shape differences;
however, its maximum deviation is dominated by the bulk of the distribution and is
relatively insensitive to tail discrepancies.

\textbf{Wasserstein Distance} \quad The Wasserstein distance measures the minimal effort
required to transform one distribution into another. For two cumulative distributions
\(F(x)\) and \(G(x)\), it is defined as:
\[
    W_1(F, G) = \int_{-\infty}^{\infty} |F(x) - G(x)|\, dx
\]
Unlike the KS statistic, which reports only the maximum deviation, \(W_1\) integrates
the discrepancy across the entire support, providing a global measure of distributional
mismatch.

\section{Aggregate Statistics}

To obtain a more comprehensive view of generative quality, the pathwise mean, variance,
skewness, and kurtosis are computed for each window in \(\mathcal{D}_{\text{train}}\) and
\(\mathcal{D}_{\text{gen}}\). Scalar metrics such as KS and \(W_1\) flatten temporal
structure; pathwise higher-order moments instead provide direct evidence of whether the
model reproduces non-Gaussian behaviour and the stylized facts of financial time series,
with particular attention to skewness and excess kurtosis. The latter must be interpreted with caution, since, as discussed in Sec.\ \ref{sec:fts_literature}, the fourth moment is likely unbounded for FTS. Hence, even if computationally finite on a limited dataset, it may be infinite for the true underlying population.

\section{Visualizations}
Complementary with the statistics we plotted different sets of graphs.

\textbf{Quantile--Quantile (QQ) Plots} \quad Quantile comparison between pooled generated and real returns are useful to check whether a model who can match the bulk if it still misses the tails. By plotting quantiles of \(D_{\text{gen}}\)\ again the quantiles of \(D_{\text{train}}\).

\textbf{ECDF Plots} \quad It is a visual complement to the KS test. Plot the Empirical Cumulative Distribution Functions \(\hat{F}_{\text{gen}}\) and \(\hat{F}_{\text{data}}\) together with the KS gap marked. This makes the KS statistic more interpretable.

\textbf{Unnormalized Paths} \quad Invert the normalisation to recover approximate price-level trajectories. This serves primarily as a sanity and plausibility check of the model's ability to generate coherent price levels, where ``unnormalize'' denotes applying the inverse of the global z-score.

\section{Financial Time Series Evaluation}
For financial time series, a set of tools is used to compute the tail distribution,
the autocorrelation function (ACF), the leverage effect, and a power-law tail
exponent \(\alpha\). These metrics are implemented following the functions provided
in \cite{stakahashy_fingan}, which was also adopted by Kim et al.\
\cite{kim2025diffusion}, enabling a direct comparison with their results. For
consistency, the power-law coefficient is computed using the same Python package:
\texttt{powerlaw} by Alstott et al.\ \cite{alstott2014powerlaw}. The definitions of the tail distribution, \(\alpha\),
volatility clustering, and the leverage effect are provided in
Sec.\ \ref{sec:fts_literature}. Following Kim et al.\ \cite{kim2025diffusion}, the heavy-tail diagnostic and the exponent \(\alpha\) are computed on the positive (right) tail only.

Following FinGAN\ \cite{stakahashy_fingan} and Kim et al.\ \cite{kim2025diffusion}, in Phase~1 the maximum lag for the leverage effect is set to 100 and for volatility clustering to 1000. In Phase~2, the maximum lag for volatility clustering is reduced to 256 for estimation stability. The ACF at lag \(k\) is estimated over \(L - k\) observation pairs; for a window of length \(L = 512\), lags approaching \(L\) leave too few pairs to produce reliable estimates, rendering a maximum lag of 512 impractical. Setting the upper bound to \(L/2 = 256\) ensures that at least half the window's values contributes to each estimate.


\section{Conditional Evaluation}
When evaluating the conditional generation we are evaluating the quality of the conditional generation by focusing on CRPS (Continuous Rank Probability Score) and interval coverage.

\textbf{CRPS} \quad This metrics is measured as:
\[
\text{CRPS}(\hat{F}, y) = \mathbb{E}_{\hat{F}} \left[|X - y|\right] - \tfrac{1}{2}\,\mathbb{E}_{\hat{F}}\left[|X - X'|\right]
\]

where \(X,X'\) are independent draws from the predictive ensemble \(\hat{F}\) and \(y\) is the observed Close. CRPS \cite{matheson1976scoring} is a proper scoring rule: it penalises inaccuracy 
(the first term) and rewards spread (the second term lowers CRPS as ensemble 
diversity grows). It reduces to MAE when the distribution collapses to a point forecast, so it can be interpreted as a calibration-adjusted MAE, where \(0 \le \text{CRPS}\le\text{MAE}\) (proof in Sec.\ \ref{app:sec:crps_bounds}). This score is also used in CSDI \cite{tashiro2021csdi}. Nevertheless, a raw CRPS value is difficult to interpret in isolation, so it was constructed a \emph{climatology} baseline whose ensemble is formed by drawing, with replacement, from the full pool of observed Close log-returns, ignoring the OHL conditioning. A second baseline fits an OLS regression of Close on the same-day Open, High, and Low log-returns, with its ensemble formed by adding bootstrapped residuals to the OLS point forecast. Relative skill is summarised by the CRPS skill score $\text{SS}_{\text{CRPS}} = 1 - \overline{\text{CRPS}}_{\text{model}} / \overline{\text{CRPS}}_{\text{baseline}}$, where $1$ is perfect, $0$ means no improvement over the uninformative baseline.

\textbf{Interval Coverage} \quad This metrics is measure for nominal coverage levels, \(1-\alpha \in \{0.5,0.8,0.9,0.95,0.99\}\). Once defined the nominal level investigated the empirical coverage is:
\[
\text{coverage}_{1-\alpha} = \frac{1}{N} \sum_{t=1}^{L} \sum_{n=1}^{N_t} \mathbf{1}\left[y_{n,t} \in \bigl[q_{\alpha/2}(n,t),\; q_{1-\alpha/2}(n,t)\bigr]\right]
\]
where \(q\) are ensemble quantiles, which means they are quantiles computed from the set of multiple forecasts the model produces for the same target value, or set of conditioning information. A well-calibrated model has the empirical value approximately identical to nominal. Positive gap means underconfidence (intervals too wide), negative gap means overconfidence (intervals too narrow).
